% GENERATED FILE. Edit canonical lessons, then run npm run generate:books.
% canonical-source-hash: fnv1a64:d48cb0550fca7245
% canonical-lessons: FR-C486-insomnie, FR-C486-ronfler, FR-C486-cauchemar, FR-C486-bailler, FR-C486-sieste

\chapter{Nuits agitées}
\label{ch:fr-nuits-agitees}
\hlchaptermodality{486}

\begin{chapteropening}
\textbf{By the end of this chapter:} \emph{I can describe a bad night: sleeplessness, snoring and nightmares.}

Complete the last lesson of chapter 486: I can describe a bad night: sleeplessness, snoring and nightmares.
\end{chapteropening}

\section[l'insomnie]{l'insomnie — insomnia, sleeplessness}

\label{lesson:FR-C486-insomnie}



\begin{quote}
Before the new one: say the French for convalescence, then the French for to look after.
\end{quote}

\subsection*{You'll want to know: l'insomnie}
\textbf{l'insomnie} — \textquotedblleft{}insomnia, sleeplessness\textquotedblright{}.

\begin{quote}
Quand je suis stressé, je souffre d'insomnie.
\end{quote}

\begin{quote}
\emph{When I am stressed, I suffer from insomnia.}
\end{quote}

Feminine. Often plural: \emph{avoir des insomnies}.

\begin{cousinweb}
\emph{l'insomnie} is a learned borrowing from Latin \emph{insomnia} \textquotedblleft{}sleeplessness\textquotedblright{} (EtymDB), built on \emph{somnus} \textquotedblleft{}sleep\textquotedblright{} (Lewis \& Short).
\end{cousinweb}

\begin{grammarlens}[title={What you've built}]
One more word for telling, explaining and describing at B1.
\end{grammarlens}

\subsection*{Your turn}
\begin{itemize}
\raggedright
  \item \emph{Say it:} \emph{l'insomnie}
  \item \emph{Say it:} \emph{l'insomnie}, once more
  \item \emph{Say it:} a sentence of your own with \emph{l'insomnie}
  \item \emph{Recall:} say the French for convalescence, then the French for to look after, then say \emph{l'insomnie} again
\end{itemize}

\begin{tcolorbox}[breakable,colback=teal!4,colframe=teal!35!black,title={Before you move on}]
What does l'insomnie mean? (\textquotedblleft{}Insomnia, sleeplessness\textquotedblright{}.) Say it once more.
\end{tcolorbox}
\section[ronfler]{ronfler — to snore}

\label{lesson:FR-C486-ronfler}



\begin{quote}
Before the new one: say the French for to look after, then the French for insomnia.
\end{quote}

\subsection*{You'll want to know: ronfler}
\textbf{ronfler} — \textquotedblleft{}to snore\textquotedblright{}.

\begin{quote}
Mon colocataire ronfle très fort.
\end{quote}

\begin{quote}
\emph{My flatmate snores very loudly.}
\end{quote}

Also of an engine: \emph{le moteur ronfle}.

\begin{cousinweb}
\emph{ronfler} goes back to Old French \emph{ronfler} (EtymDB); it is usually taken to imitate the sound, perhaps with help from \emph{souffler} \textquotedblleft{}to blow\textquotedblright{}.
\end{cousinweb}

\begin{grammarlens}[title={What you've built}]
One more word for telling, explaining and describing at B1.
\end{grammarlens}

\subsection*{Your turn}
\begin{itemize}
\raggedright
  \item \emph{Say it:} \emph{ronfler}
  \item \emph{Say it:} \emph{ronfler}, once more
  \item \emph{Say it:} a sentence of your own with \emph{ronfler}
  \item \emph{Recall:} say the French for to look after, then the French for insomnia, then say \emph{ronfler} again
\end{itemize}

\begin{tcolorbox}[breakable,colback=teal!4,colframe=teal!35!black,title={Before you move on}]
What does ronfler mean? (\textquotedblleft{}To snore\textquotedblright{}.) Say it once more.
\end{tcolorbox}
\section[le cauchemar]{le cauchemar — nightmare}

\label{lesson:FR-C486-cauchemar}



\begin{quote}
Before the new one: say the French for insomnia, then the French for to snore.
\end{quote}

\subsection*{You'll want to know: le cauchemar}
\textbf{le cauchemar} — \textquotedblleft{}nightmare\textquotedblright{}.

\begin{quote}
Cette nuit, j'ai fait un cauchemar terrible.
\end{quote}

\begin{quote}
\emph{Last night I had a terrible nightmare.}
\end{quote}

\emph{Faire un cauchemar} — to have a nightmare.

\begin{cousinweb}
\emph{le cauchemar} comes from Old French \emph{cauquemare} (EtymDB), built from \emph{caucher} \textquotedblleft{}to press\textquotedblright{}, from Latin \emph{calcāre} \textquotedblleft{}to tread\textquotedblright{}, and Middle Dutch \emph{mare} \textquotedblleft{}a night ghost\textquotedblright{}, a Germanic word for a spirit that sits on the sleeper.
\end{cousinweb}

\begin{grammarlens}[title={What you've built}]
One more word for telling, explaining and describing at B1.
\end{grammarlens}

\subsection*{Your turn}
\begin{itemize}
\raggedright
  \item \emph{Say it:} \emph{le cauchemar}
  \item \emph{Say it:} \emph{le cauchemar}, once more
  \item \emph{Say it:} a sentence of your own with \emph{le cauchemar}
  \item \emph{Recall:} say the French for insomnia, then the French for to snore, then say \emph{le cauchemar} again
\end{itemize}

\begin{tcolorbox}[breakable,colback=teal!4,colframe=teal!35!black,title={Before you move on}]
What does le cauchemar mean? (\textquotedblleft{}Nightmare\textquotedblright{}.) Say it once more.
\end{tcolorbox}
\section[bâiller]{bâiller — to yawn}

\label{lesson:FR-C486-bailler}



\begin{quote}
Before the new one: say the French for to snore, then the French for nightmare.
\end{quote}

\subsection*{You'll want to know: bâiller}
\textbf{bâiller} — \textquotedblleft{}to yawn\textquotedblright{}.

\begin{quote}
Les élèves bâillent pendant le cours de l'après-midi.
\end{quote}

\begin{quote}
\emph{The pupils yawn during the afternoon lesson.}
\end{quote}

Do not mix it up with \emph{bailler}, an old word for giving.

\begin{cousinweb}
\emph{bâiller} is inherited from Old French \emph{baaillier}, from Vulgar Latin \emph{bataclāre} (EtymDB), built on \emph{batāre} \textquotedblleft{}to gape\textquotedblright{}.
\end{cousinweb}

\begin{grammarlens}[title={What you've built}]
One more word for telling, explaining and describing at B1.
\end{grammarlens}

\subsection*{Your turn}
\begin{itemize}
\raggedright
  \item \emph{Say it:} \emph{bâiller}
  \item \emph{Say it:} \emph{bâiller}, once more
  \item \emph{Say it:} a sentence of your own with \emph{bâiller}
  \item \emph{Recall:} say the French for to snore, then the French for nightmare, then say \emph{bâiller} again
\end{itemize}

\begin{tcolorbox}[breakable,colback=teal!4,colframe=teal!35!black,title={Before you move on}]
What does bâiller mean? (\textquotedblleft{}To yawn\textquotedblright{}.) Say it once more.
\end{tcolorbox}
\section[la sieste]{la sieste — nap, siesta}

\label{lesson:FR-C486-sieste}



\begin{quote}
Before the new one: say the French for nightmare, then the French for to yawn.
\end{quote}

\subsection*{You'll want to know: la sieste}
\textbf{la sieste} — \textquotedblleft{}nap, siesta\textquotedblright{}.

\begin{quote}
Après le déjeuner, mon grand-père fait la sieste.
\end{quote}

\begin{quote}
\emph{After lunch, my grandfather has a nap.}
\end{quote}

\emph{Faire la sieste} — to have a nap.

\begin{cousinweb}
\emph{la sieste} was borrowed from Spanish \emph{siesta} (EtymDB), from Latin \emph{sexta} (\emph{hōra}) \textquotedblleft{}the sixth hour\textquotedblright{}, that is noon.
\end{cousinweb}

\begin{grammarlens}[title={What you've built}]
One more word for telling, explaining and describing at B1.
\end{grammarlens}

\subsection*{Your turn}
\begin{itemize}
\raggedright
  \item \emph{Say it:} \emph{la sieste}
  \item \emph{Say it:} \emph{la sieste}, once more
  \item \emph{Say it:} a sentence of your own with \emph{la sieste}
  \item \emph{Recall:} say the French for nightmare, then the French for to yawn, then say \emph{la sieste} again
\end{itemize}

\begin{tcolorbox}[breakable,colback=teal!4,colframe=teal!35!black,title={Before you move on}]
What does la sieste mean? (\textquotedblleft{}Nap, siesta\textquotedblright{}.) Say it once more.
\end{tcolorbox}
